% Copy-ready AMS-style front matter for the Problem 593 manuscript.
% This file is not included automatically.  It preserves the mathematical
% claims while normalising title metadata, abstract structure, theorem prose,
% attribution, and the proof roadmap.

\title[Obligatory triple systems]
  {Obligatory Triple Systems: Bridge Decompositions and Exact Finite Spectra}
\author{Samuil Petkov}
\address{École normale supérieure--PSL, 45 rue d'Ulm, 75005 Paris, France}
\email{samuil.petkov@ens.psl.eu}
\subjclass[2020]{Primary 05C65; Secondary 05C15, 05C63, 03E05}
\keywords{obligatory triple system, weak hypergraph coloring, Levi graph,
  Berge cycle, uncountable chromatic number, Erdős Problem 593}

\begin{abstract}
We give an alternative proof of the classification of finite obligatory triple
systems.  After isolated vertices are removed, obligatoriness is equivalent to
linearity, the presence of a bridge incident with every hyperedge-node of the
Levi graph, and evenness of every Berge cycle.  Equivalently, the system is
obtained from private-vertex expansions of finite bipartite graphs and finite
edgeless systems by disjoint unions and one-point amalgamations.  The proof
separates a finite bridge-block decomposition from the infinitary avoiding-host
construction based on a one-apex sequence lift.  We also determine the exact
order--size--component spectrum: a reduced obligatory system with $m\ge1$
hyperedges and $c$ connected components has $n$ vertices exactly when
$m+2(c-1)+\lceil2\sqrt{m-c+1}\rceil\le n\le2m+c$.  A Lean~4 development verifies
the finite classification.
\end{abstract}

\maketitle

\section{Introduction}\label{introduction-ams}

The hosts in Erdős Problem~593 are infinite, but the classification is governed
by a finite incidence structure.  A finite triple system is \emph{obligatory}
if it embeds in every triple system of uncountable weak chromatic number.
Erdős asked for a characterization of these systems~\citep{erdos1995}; the
problem is catalogued as Erdős Problem~593~\citep{bloom593}.

For a finite graph $J$, let $J^+$ denote its private-vertex expansion: every
edge $uv$ is replaced by a triple $\{u,v,p_{uv}\}$, where the points $p_{uv}$
are new and pairwise distinct.  Let $\mathcal B$ be the smallest class that
contains $J^+$ for every finite bipartite graph $J$, contains every finite
edgeless triple system, and is closed under finite disjoint unions and
one-point amalgamations.  Membership in $\mathcal B$ is understood up to
triple-system isomorphism.  For a triple system $F$, let $F^\circ$ be the system
obtained by deleting all isolated vertices, and let $I(F)$ denote its Levi, or
incidence, graph.

\begin{theoremA}\label{theorem-a-ams}
For every finite triple system $F$, the following assertions are equivalent.
\begin{enumerate}
\item $F$ is obligatory.
\item $F\in\mathcal B$.
\item The system $F^\circ$ is linear, every hyperedge-node of $I(F^\circ)$ is
incident with a bridge, and every Berge cycle of $F^\circ$ has even length.
\end{enumerate}
\end{theoremA}

Li's preprint contains the first publicly posted complete proof of this
classification and introduces the complete-rank one-apex lift and bridge-trace
method used in the negative direction
\citep[Theorem~1.1 and Sections~3--4]{li2026}.  The present paper gives an
alternative implementation, a direct all-bridges decomposition of the finite
structure, exact finite parameter consequences, and a Lean formalization of
this implementation.  No competing priority claim is made for the
classification or the shared lift architecture.  The positive expansion atoms
also follow from Reiher's theorem on obligatory hypergraphs
\citep[Theorem~1.2]{reiher2024}; we include a direct proof in the present
notation.

The proof separates into three implications:
\[
\mathcal B\Longrightarrow\text{obligatory},\qquad
\mathcal B\Longleftrightarrow\text{intrinsic},\qquad
\neg\text{intrinsic}\Longrightarrow\neg\text{obligatory}.
\]
The middle equivalence is the finite geometric core.  Delete all bridges from
the Levi graph.  Every hyperedge-node then has residual degree zero or two.
Each active component suppresses to a simple bipartite graph, and the deleted
incidences become distinct private points.  Contracting the bridge-deleted
components produces a forest, whose rooted order reconstructs the triple system
by disjoint unions and one-point amalgamations.

For the positive implication, complete bipartite expansions are obligatory, and
obligatoriness is preserved by finite disjoint unions and one-point
amalgamations.  The direct proof of the expansion atom uses a locally bounded
rainbow lemma, while the amalgamation argument uses rooted abundance.

For necessity, the intrinsic condition can fail in exactly three ways.
A classical Erdős--Rado construction supplies an uncountably chromatic linear
triple system and excludes every nonlinear source.  A one-apex lift of
$K_{\omega_1}$ excludes a linear source whose Levi graph has a hyperedge-node
without an incident bridge.  Finally, lifting an uncountably chromatic graph
whose odd girth exceeds a prescribed finite bound excludes a source containing
an odd Berge cycle.

The bridge-block proof also retains exact numerical information.  It associates
to every reduced obligatory system $F$ a finite bipartite shadow $J$ with the
same numbers of hyperedges and connected components and with
\[
|V(J)|=|V(F)|-|E(F)|.
\]
Ordinary extremal bounds for finite bipartite graphs then give the complete
order--size--component spectrum and its fixed-order and cycle-rank corollaries.

The paper is organized as follows.  The preliminaries fix the coloring,
embedding, incidence, and cycle conventions and record the external inputs.
We then prove the finite bridge-block theorem, the positive implication, the
one-apex trace theorem, and the three avoiding-host propositions.  The final
mathematical section derives the exact finite spectra.  The last section states
the formal verification and reproducibility boundary.